% !PRINTED 526
\weberpage{526}
%<<CONT
(durch Adjunction von Variablen) vergrössert wird. Es entsteht so ein erweiterter Körper, in dem die Gesetze der Zerlegbarkeit rein gelten.

Den Ausgangspunkt der Definitionen bildete der Körper $R$ der rationalen Zahlen, und wenn wir uns Rechenschaft darüber geben wollen, auf welchen Eigenschaften des Körpers $R$ die Möglichkeit dieser Erweiterung beruht, so finden wir, dass es einerseits die Existenz der ganzen Zahlen in $R$, andererseits die eindeutige Zerlegbarkeit dieser ganzen Zahlen in Primfactoren ist, die allein bei der ganzen Deduction benutzt wurden. Wenn wir also an Stelle des Körpers $R$ irgend einen anderen Körper treten lassen, dem diese beiden Eigenschaften zukommen, so werden wir dieselben Folgerungen ziehen können. Nehmen wir für $R$ einen anderen algebraischen Zahlkörper, so bekommen wir freilich nichts Neues, wohl aber, wenn wir z.\,B. an Stelle des Körpers $R$ den Körper der rationalen Functionen einer Variablen setzen, dem ja die beiden fundamentalen Eigenschaften auch zukommen. So gewinnen wir einen Ausgangspunkt für die Theorie der algebraischen Functionen einer Variablen\footnote{Vergl. \gesperrt{Dedekind-Weber}, Theorie der algebraischen Functionen einer Veränderlichen. \gesperrt{Crelle's} Journal, Bd. 92.}.

\begin{center}
\rule{4cm}{0.4pt}
\end{center}
